%% Resumo
\begin{resumo}
Este trabalho apresenta o projeto desenvolvido para a disciplina de Projeto Integrador de Extensão I do curso de Análise e Desenvolvimento de Sistemas do Instituto Federal de Educação, Ciência e Tecnologia de São Paulo. A proposta consiste no desenvolvimento de um sistema web para apoio à gestão do Sebo do Cesar, estabelecimento especializado na comercialização de livros usados. O projeto foi elaborado com base nos conhecimentos adquiridos ao longo do curso e fundamentado em práticas de engenharia de software, contemplando levantamento de requisitos, regras de negócio, modelagem de dados, definição da arquitetura da solução e seleção das tecnologias de desenvolvimento. Além de atender às necessidades do parceiro, o projeto possibilita a aplicação prática dos conhecimentos acadêmicos em um contexto real, fortalecendo a integração entre ensino, tecnologia e extensão.
\vspace{\onelineskip}
\noindent

\textbf{Palavras-chaves}: Gestão; Sebo; Livros; Sistema Web; Desenvolvimento de Software.
\end{resumo}